% Bewertungspaket zum Gesamttest «Lineare und quadratische Gleichungen» — Quelle des PDFs (python3 scripts/build-lp-pdf.py lineare-quadratische).
% Ganze Punkte je Teilschritt; Werte mit python3 nachgerechnet (06.10.2026, Folgefehler-Fälle eingeschlossen).
\documentclass[11pt]{article}
\usepackage{../lp-druck}
\renewcommand{\lpname}{Lineare und quadratische Gleichungen}
\renewcommand{\lpbereich}{Grundlagenfach}
\renewcommand{\lpteilgebiet}{2.2}
\renewcommand{\lpfuss}{Bewertungspaket}
\newcolumntype{P}{>{\centering\arraybackslash}p{10mm}}
\newenvironment{raster}{\par\smallskip\noindent\tabularx{\linewidth}{@{}>{\raggedright\arraybackslash}p{38mm}P X@{}}\toprule
  \small\textsc{Teilschritt} & \small\textsc{P} & \small\textsc{Musterlösung}\\\midrule}{\bottomrule\endtabularx\par}
\newcommand{\fehler}[1]{\par\smallskip{\small\textcolor{lprot}{\bfseries Typische Fehler:} #1}}
\newcommand{\E}{\,{\footnotesize\bfseries\color{lpgruen}(E)}}
\newcommand{\A}{\,{\footnotesize\bfseries\color{lpblau}(A)}}
\begin{document}
\lpkopf{Bewertungspaket zum Gesamttest}{}

\begin{lpachtung}{Erst lösen, dann öffnen}
Dieses Paket enthält die Musterlösung. Wer es vor dem Lösen liest, misst am Ende nur, wie gut das Abschreiben gelingt.
\end{lpachtung}

\section*{1 · So gehst du vor}
\begin{enumerate}[leftmargin=*,itemsep=2pt]
\item \textbf{Gesamttest lösen} — vollständig, mit Rechenweg, ohne dieses Paket und ohne Taschenrechner.
\item \textbf{Fotografieren oder scannen}, gut lesbar, eine Seite pro Bild. \textbf{Kein Name} auf den Bildern.
  Handyfotos tragen oft unsichtbar Standort, Zeit und Gerät mit (Metadaten). Lade darum lieber einen Scan oder ein
  Bildschirmfoto des Fotos hoch, oder entferne vorher die Standortangaben.
\item \textbf{Bewerten lassen.} Ob du dafür eine KI nutzt und welche, entscheidest du selbst und in eigener Verantwortung —
  je nachdem, was dir aufgrund deines Alters und deiner persönlichen Situation erlaubt ist (Altersgrenzen und
  Nutzungsbedingungen des Anbieters, allenfalls Einverständnis deiner Eltern). Lade nur deine Lösung und dieses PDF hoch,
  keine persönlichen Daten, und füge den Auftrag aus Abschnitt~2 ein. Ohne KI kannst du dich mit Abschnitt~3 selbst bewerten.
\item \textbf{Rückmeldung prüfen.} Stimmt, was die KI aus deiner Handschrift gelesen hat? Eine KI kann sich irren —
  im Zweifel gilt das Raster in Abschnitt~3.
\item \textbf{Gezielt wiederholen} nach Abschnitt~4.
\end{enumerate}

\section*{2 · Auftrag an die KI}
\begin{tcolorbox}[colback=lplinie!25,colframe=lplinie,boxrule=0.5pt,arc=1mm,fontupper=\small]
Du bist Korrektorin oder Korrektor für Mathematik an einer Schweizer Berufsmaturitätsschule. Im Anhang findest du (1)~das Bewertungspaket zum Gesamttest «Lineare und quadratische Gleichungen» mit Musterlösung und Punkteraster und (2)~Fotos meiner handschriftlichen Lösung.\par\medskip
Geh so vor:\par
1. Schreib zu jeder Aufgabe G1 bis G8 zuerst ab, was du in meiner Lösung liest — Rechenweg und Ergebnis. Was du nicht sicher lesen kannst, markierst du mit [unsicher]. Nicht raten.\par
2. Vergib die Punkte genau nach dem Raster im Bewertungspaket (Abschnitt 3), ganze Punkte je Zeile. Ziehe einen Fehler nur einmal ab: Rechne ich mit einem falschen Zwischenergebnis richtig weiter, gibt es die Folgepunkte. Ausnahme: Zeilen mit \textbf{(E)} gibt es nur für das richtige Ergebnis, nie als Folgepunkt. Die Zeilen «Typische Fehler» sagen, welche Rasterzeile bei diesem Fehler 0 Punkte gibt — sie sind kein zusätzlicher Abzug.\par
3. Andere richtige Lösungswege zählen voll. Gleichwertige Schreibweisen zählen voll: Dezimalpunkt oder Dezimalkomma, Bruch oder Dezimalzahl ($\tfrac12 = 0.5$), gekürzter oder ungekürzter Bruch, $\mathbb{L} = \{-2;\, 4\}$ oder «$x = -2$ oder $x = 4$» oder $x_{1,2} = \ldots$, die Elemente von $\mathbb{L}$ in beliebiger Reihenfolge, Strichpunkt zwischen ihnen (bei Dezimalpunkt auch Komma), $\mathbb{L} = \{\,\}$ oder «keine Lösung», $\mathbb{L} = \mathbb{R}$ oder «jede Zahl ist Lösung» — ausser wo das Raster eine bestimmte Form verlangt. Zeilen mit \textbf{(E)} sind Ergebnispunkte: Das Ergebnis muss stimmen \emph{und} aus einem erkennbaren Weg hervorgehen, ausser die Zeile sagt ausdrücklich etwas anderes. Alle übrigen Zeilen bewerten den Weg selbst.\par
4. Begründe jeden Abzug in einem Satz und nenne die Fehlerart (zum Beispiel «Minus vor der Klammer» oder «Fall $k = 1$ übersehen»). Schreib die Musterlösung nicht ab — zeig mir, wo mein Fehler liegt.\par
5. Gib zum Schluss eine Tabelle «Aufgabe | Punkte | Begründung», die Gesamtpunktzahl (von 25) und — nach der Selbsteinschätzung in Abschnitt 4 — welches Kapitel des Leitprogramms ich wiederholen soll. Keine Note.\par
6. Fehlt ein Foto oder ist eine Aufgabe unlesbar, sag es, statt sie zu bewerten.
\end{tcolorbox}

\needspace{16\baselineskip}
\section*{3 · Musterlösung und Punkteraster}
\textbf{Allgemein:} Ganze Punkte je Zeile. Ein Fehler wird nur einmal abgezogen (Folgefehler). Gleichwertige Wege und
Schreibweisen zählen voll. In der Spalte P: \textbf{(E)}~= Ergebnispunkt, Ergebnis richtig \emph{und} Weg erkennbar
(ausser die Zeile sagt ausdrücklich etwas anderes); ohne Zeichen = die Zeile bewertet den Weg oder die Begründung.
Der ganze Test ist ohne Taschenrechner gestellt; exakte Werte (Brüche, Wurzeln) sind verlangt, wo die Aufgabe es sagt.
«Typische Fehler» nennt, welche Zeile 0 Punkte gibt, und die Punkte, die bleiben.
\textbf{Teilrichtige Zeilen:} Verlangt eine Zeile mehrere Angaben und stimmt nur ein Teil, gibt die Zeile 0 Punkte —
ausser die Zeile nennt ausdrücklich eine Aufteilung.
\textbf{Komma:} Wer das Dezimalkomma benutzt, muss die Elemente von $\mathbb{L}$ mit Strichpunkt trennen. Steht «$\{0,5\}$»
ohne Strichpunkt und ist aus dem Weg nicht klar, ob $0.5$ oder $0$ und $5$ gemeint ist: [unsicher] melden, nicht raten.

\begin{lpaufgabe}{G1}{Lineare Gleichung mit Brüchen}{3 P}
\begin{raster}
Hauptnenner und Klammer & 1 & $\cdot 4$: $x + 2 - 4(x - 3) = 2(x - 1)$, also $x + 2 - 4x + 12 = 2x - 2$, d.\,h. $-3x + 14 = 2x - 2$. Jedes Glied mal $4$, das Minus vor der Klammer wirkt auf beide Glieder.\\
Lösung & 1\E & $16 = 5x$, $x = \tfrac{16}{5} = 3.2$, $\mathbb{L} = \{3.2\}$\\
Probe & 1 & links $\tfrac{5.2}{4} - 0.2 = 1.3 - 0.2 = 1.1$, rechts $\tfrac{2.2}{2} = 1.1$ — in die \emph{Ausgangsgleichung} eingesetzt. Gilt auch als Folgefehler: Eine richtig gerechnete Probe zur eigenen falschen Lösung, die den Fehler aufdeckt (beide Seiten verschieden, so benannt), gibt den Punkt.\\
\end{raster}
\fehler{$-(x - 3)$ als $-x - 3$: $-3x - 10 = 2x - 2$, $x = -1.6$ — erste und Lösungszeile 0, Probe nach obiger Regel $\to$ höchstens 1 von 3\,P.
Nur die Brüche mal $4$ genommen ($x + 2 - (x - 3) = 2(x - 1)$, $x = 3.5$): erste und Lösungszeile 0.
Probe in einer umgeformten Zeile statt in der Ausgangsgleichung: Probezeile 0.}
\end{lpaufgabe}

\begin{lpaufgabe}{G2}{Lösungsfälle}{2 P}
\begin{raster}
Umformen & 1 & $3x - 6 = 3x + c$, nach $\mid -3x$: $-6 = c$ — das $x$ fällt weg, es bleibt eine Aussage über $c$.\\
Fälle & 1\E & $c = -6$: $-6 = -6$ wahr, jede Zahl ist Lösung, $\mathbb{L} = \mathbb{R}$. Jedes andere $c$: falsche Aussage, $\mathbb{L} = \{\,\}$. Beide Fälle nötig.\\
\end{raster}
\fehler{Nur Beispiele ($c = -6$ und z.\,B. $c = 0$) ohne die Aussage, die übrig bleibt: Umformzeile 0, Fälle-Zeile 1.
$3(x - 2) = 3x - 2$ ausmultipliziert, also $c = -2$: Umformzeile 0; Fälle mit $-2$ richtig weitergedacht: Fälle-Zeile 0 (E).
«Für $c = -6$ keine Lösung»: Fälle-Zeile 0.}
\end{lpaufgabe}

\begin{lpaufgabe}{G3}{Nullprodukt — oder nicht?}{3 P}
\begin{raster}
(a) & 1\E & $x + 2 = 0$ oder $x - 3 = 0$: $\mathbb{L} = \{-2;\, 3\}$. Beide nötig.\\
(b) & 1\E & Ausmultiplizieren: $x^2 - x - 6 = 14$, auf null: $x^2 - x - 20 = 0$, $(x - 5)(x + 4) = 0$ (oder Mitternachtsformel, $D = 81$): $\mathbb{L} = \{-4;\, 5\}$\\
(c) & 1 & Der Satz vom Nullprodukt gilt nur, wenn das Produkt \emph{null} ist. Ist es $14$, muss kein Faktor einen bestimmten Wert haben (z.\,B. $2 \cdot 7 = 14$). Gleichwertig: Probe zeigt, dass $x = 12$ bzw. $x = 17$ falsch sind.\\
\end{raster}
\fehler{(b) «$x + 2 = 14$ oder $x - 3 = 14$», $\mathbb{L} = \{12;\, 17\}$: Zeile (b) 0; (c) kann trotzdem stimmen.
(b) Ausmultipliziert, aber nicht auf null gebracht ($x^2 - x - 6 = 14$ direkt «gelöst»): Zeile (b) 0.
(b) Vorzeichen beim Faktorisieren, $\mathbb{L} = \{-5;\, 4\}$: Zeile (b) 0.}
\end{lpaufgabe}

\begin{lpaufgabe}{G4}{Quadratische Ergänzung mit Zahl \texorpdfstring{$c$}{c}}{3 P}
\begin{raster}
(a) Ergänzung & 1 & $x^2 - 6x = -c$; beidseitig $+\left(\tfrac{6}{2}\right)^2 = +9$: $(x - 3)^2 = 9 - c$\\
(b) Anzahl & 1 & $9 - c > 0$, also $c < 9$: zwei Lösungen; $c = 9$: eine ($x = 3$); $c > 9$: keine — ein Quadrat ist nie negativ. Alle drei Fälle nötig.\\
(c) & 1\E & $c = 5$: $(x - 3)^2 = 4$, $x - 3 = \pm 2$, $\mathbb{L} = \{1;\, 5\}$\\
\end{raster}
\fehler{(a) $+9$ nur links addiert: $(x - 3)^2 = -c$ — Ergänzungszeile 0; (b) mit dieser Form richtig weitergedacht ($c < 0$ zwei usw.): Zeile (b) als Folgefehler 1; (c) dann falsch: 0 $\to$ 1 von 3\,P.
(b) über die Diskriminante $D = 36 - 4c$ statt über die Ergänzung: gleichwertig, Zeile (b) voll.
(c) Nur $+2$ beim Wurzelziehen, $\mathbb{L} = \{5\}$: Zeile (c) 0.}
\end{lpaufgabe}

\begin{lpaufgabe}{G5}{Mitternachtsformel, exakt}{3 P}
\begin{raster}
$D$ & 1 & $a = 1$, $b = -2$, $c = -1$; $D = 4 + 4 = 8$\\
Einsetzen & 1 & $x_{1,2} = \dfrac{2 \pm \sqrt8}{2}$ (Vorzeichen von $-b$ beachtet)\\
Lösung & 1\E & $x = 1 \pm \sqrt2$, $\mathbb{L} = \{1 - \sqrt2;\, 1 + \sqrt2\}$. Gleichwertig: $\tfrac{2 \pm \sqrt8}{2}$ oder $\tfrac{2 \pm 2\sqrt2}{2}$. Nur gerundet ($\approx -0.41$; $2.41$) ohne exakte Form: 0 — die Aufgabe verlangt exakt.\\
\end{raster}
\fehler{$D = 4 - 4 = 0$ (Vorzeichen von $-4ac$ bei $c = -1$): $D$-Zeile 0; mit $D = 0$ richtig weitergerechnet ($x = 1$): Einsetzzeile 1, Lösungszeile 0.
$-b = -2$ eingesetzt: $\tfrac{-2 \pm \sqrt8}{2} = -1 \pm \sqrt2$ — Einsetz- und Lösungszeile 0.
Nur durch $2$ statt durch $2a$ geteilt: hier gleich ($a = 1$), kein Abzug.}
\end{lpaufgabe}

\begin{lpaufgabe}{G6}{Verfahren begründen und prüfen}{5 P}
\begin{raster}
(a) & 1 & Wurzelziehen: kein lineares Glied ($b = 0$), nach $x^2$ auflösen. Gleichwertig: Binom $(2x - 3)(2x + 3) = 0$ (Differenz zweier Quadrate). Verfahren \emph{und} Grund nötig.\\
(b) & 1 & Binom erkennen: $x^2 + 6x + 9 = (x + 3)^2$. Gleichwertig: Zweiklammersatz mit $3$ und $3$. Verfahren und Grund nötig.\\
(c) & 1 & Mitternachtsformel: Faktor $2$ vor $x^2$, alle drei Glieder da, keine Zerlegung mit ganzen Zahlen erkennbar ($D = 33$). Verfahren und Grund nötig.\\
(d) Probe & 1 & $x = 2$: $4 + 2 - 6 = 0$ stimmt; $x = 3$: $9 + 3 - 6 = 6 \neq 0$ — falsch. Beide Proben nötig.\\
(d) Korrektur & 1\E & $(x - 2)(x + 3) = 0$ (Produkt $-6$, Summe der Lösungen $-1$): richtig ist $\mathbb{L} = \{-3;\, 2\}$\\
\end{raster}
\fehler{Verfahren ohne Begründung genannt: Zeile 0. Mitternachtsformel bei (a) oder (b) mit Grund «geht immer»: Zeile 0 — gefragt ist das schnellste Verfahren.
(d) Nur «3 ist falsch» ohne eingesetzte Probe: Probezeile 0. Statt $-3$ eine andere Zahl ohne Probe angegeben: Korrekturzeile 0.}
\end{lpaufgabe}

\begin{lpaufgabe}{G7}{Lineare Gleichung mit Parameter}{3 P}
\begin{raster}
Kritische Werte & 1 & Der Faktor vor $x$ ist $k(k - 1)$; er wird null für $k = 0$ und $k = 1$. Beide Werte nötig.\\
Sonderfälle & 1 & $k = 0$: $0 \cdot x = 0$, $\mathbb{L} = \mathbb{R}$. $k = 1$: $0 \cdot x = 1$, $\mathbb{L} = \{\,\}$. Beide nötig.\\
Allgemeiner Fall & 1\E & $k \neq 0$, $k \neq 1$: $x = \dfrac{k}{k(k - 1)} = \dfrac{1}{k - 1}$, $\mathbb{L} = \left\{\tfrac{1}{k - 1}\right\}$. Ungekürzt $\tfrac{k}{k(k - 1)}$ zählt voll.\\
\end{raster}
\fehler{Nur einen kritischen Wert gefunden (meist $k = 1$) und diesen Fall richtig behandelt: erste Zeile 0, Sonderfälle-Zeile 0, allgemeiner Fall voll, wenn «$k \neq 1$» und $\tfrac{1}{k - 1}$ dastehen $\to$ 1 von 3\,P.
Ohne Fallunterscheidung durch $k(k - 1)$ geteilt: erste Zeile 0, Sonderfälle 0, allgemeiner Fall 1.
Bei $k = 0$ «keine Lösung» oder bei $k = 1$ «alle»: Sonderfälle-Zeile 0.}
\end{lpaufgabe}

\begin{lpaufgabe}{G8}{Quadratische Gleichung mit Parameter}{3 P}
\begin{raster}
Fall $k = 1$ & 1 & Faktor vor $x^2$ null: $2x + 1 = 0$, linear, genau eine Lösung $x = -0.5$\\
$D(k)$ & 1 & $k \neq 1$: $D = 2^2 - 4(k - 1) \cdot 1 = 8 - 4k$\\
Fälle & 1\E & $k = 2$: $D = 0$, eine Lösung $x = -\tfrac{2}{2 \cdot 1} = -1$. $k < 2$ (und $k \neq 1$): zwei Lösungen. $k > 2$: keine. Alle drei nötig.\\
\end{raster}
\fehler{Fall $k = 1$ übersehen, nur mit $D$ gerechnet: erste Zeile 0, die übrigen voll, wenn sie stimmen $\to$ 2 von 3\,P (beim Fälle-Punkt reicht «$k < 2$» ohne «$k \neq 1$» dann nicht: 1 von 3\,P).
$D = 4 - 4k$ (die $-1$ in $k - 1$ vergessen): $D$-Zeile 0; Fälle mit diesem $D$ ($k = 1$ eine Lösung …): Fälle-Zeile 0 (E).
Bei $k = 2$ keine Lösung angegeben: Fälle-Zeile 0.}
\end{lpaufgabe}


\section*{4 · Selbsteinschätzung}
\begin{tabularx}{\linewidth}{@{}l X@{}}\toprule
22 – 25 P & Die geprüften Teile sitzen. Wo du Punkte verloren hast: das Kapitel dieser Aufgabe nochmals (Zuordnung unten).\\
17 – 21 P & Den schwächsten Teil nochmals: Tüfteln und Übungen des Kapitels, in dem du die meisten Punkte verloren hast.\\
11 – 16 P & Zurück zu den Kapiteln aller Aufgaben, in denen du Punkte verloren hast.\\
0 – 10 P & Zurück zu Kapitel 1 und von dort der Reihe nach weiter.\\\midrule
\multicolumn{2}{@{}p{\linewidth}@{}}{\small\textbf{Aufgabe $\to$ Kapitel:} G1, G2 $\to$ Kapitel 1 (Lineare Gleichungen umformen); G3 $\to$ Kapitel 2 (Ausklammern und Nullprodukt); G4, G5 $\to$ Kapitel 3 (Wurzelziehen, Ergänzen, Mitternachtsformel); G6 $\to$ Kapitel 4 (Das passende Verfahren); G7, G8 $\to$ Kapitel 5 (Parameterdiskussion)}\\\bottomrule
\end{tabularx}
\end{document}
